% Part: incompleteness
% Chapter: arithmetization-syntax
% Section: proofs-in-lk

\documentclass[../../../include/open-logic-section]{subfiles}

\begin{document}

\olfileid{inc}{art}{plk}
\olsection{$\Log{LK}$ मधील \usetoken{P}{derivation}}

\begin{explain}
!!{derivation}s चे अंकगणितीकरण करण्यासाठी !!{derivation}s संख्यांनी दर्शवावे लागतात.
प्रत्येक निगमनाला नामांकन असलेले क्रमवर्तींचे वृक्ष म्हणजे !!{derivation}s
असल्यामुळे, पुनरावर्ती निरूपण हा सहज मार्ग आहे. आपण !!{derivation}
एका क्रमित समूहाने दर्शवतो: त्याचे घटक म्हणजे अंतिम क्रमवर्ती, नामांकन,
आणि अखेरच्या निगमनाच्या आधारविधानांपर्यंत येणाऱ्या उप-!!{derivation}s ची निरूपणे.
\end{explain}

\begin{defn}
$\Gamma$ ही !!{sentence}s ची सांत क्रमिका, $\Gamma =
\tuple{!A_1, \dots, !A_n}$, असेल, तर $\Gn{\Gamma} = \tuple{\Gn{!A_1},
  \dots, \Gn{!A_n}}$.

$\Gamma \Sequent \Delta$ ही क्रमवर्ती असेल, तर
$\Gamma \Sequent \Delta$ ची गोडेल संख्या
पुढीलप्रमाणे असते:
\[
\Gn{\Gamma \Sequent \Delta} = \tuple{\Gn{\Gamma}, \Gn{\Delta}}
\]

$\pi$ ही $\Log{LK}$ मधील !!a{derivation} असेल, तर $\Gn{\pi}$ ची
व्याख्या पुढीलप्रमाणे करतो:
\begin{enumerate}
\item $\pi$ मध्ये फक्त $\Gamma \Sequent \Delta$ ही आरंभीची क्रमवर्ती
  असेल, तर $\Gn{\pi}$ म्हणजे
  \[
    \tuple{0, \Gn{\Gamma \Sequent \Delta}}.
  \]
\item $\pi$ चा शेवट एक किंवा दोन आधारविधानांच्या निगमनाने होत असेल,
  त्या निगमनाचा निष्कर्ष $\Gamma \Sequent \Delta$ असेल, आणि $\pi_1$
  व $\pi_2$ या अखेरच्या निगमनाच्या आधारविधानांवर संपणाऱ्या तात्काळ
  उपनिष्पत्ती असतील, तर $\Gn{\pi}$ अनुक्रमे
  \begin{align*}
    & \tuple{1, \Gn{\pi_1}, \Gn{\Gamma \Sequent \Delta}, k} \text{ किंवा}\\
    & \tuple{2, \Gn{\pi_1}, \Gn{\pi_2}, \Gn{\Gamma \Sequent \Delta}, k},
  \end{align*}
  असते. अखेरच्या निगमनात कोणता नियम वापरला आहे त्यानुसार $k$ पुढील
  तक्त्यात दिले आहे:

  \begin{tabular}{lccccccc}
    \text{नियम:} & \LeftR{\Weakening} & \RightR{\Weakening} &
    \LeftR{\Contraction} & \RightR{\Contraction} &
    \LeftR{\Exchange} & \RightR{\Exchange} \\
    $k$: & 1 & 2 & 3 & 4 & 5 & 6 \\[2ex]
    \text{नियम:} & \LeftR{\lnot} & \RightR{\lnot} &
    \LeftR{\land} & \RightR{\land} &
    \LeftR{\lor} & \RightR{\lor} \\
    $k$: & 7 & 8 & 9 & 10 & 11 & 12 \\[2ex]
    \text{नियम:} & \LeftR{\lif} & \RightR{\lif} &
    \LeftR{\lforall} & \RightR{\lforall} &
    \LeftR{\lexists} & \RightR{\lexists} \\
    $k$: & 13 & 14 & 15 & 16 & 17 & 18 \\[2ex]
    \text{नियम:} & \Cut & = \\
    $k$: & 19 & 20
  \end{tabular}
\end{enumerate}
\end{defn}

\begin{ex}
  पुढील अगदी सोपी !!{derivation} पाहा:
  \begin{prooftree}
    \Axiom$!A \fCenter !A$
    \RightLabel{\LeftR{\land}}
    \UnaryInf$!A \land !B \fCenter !A$
    \RightLabel{\RightR{\lif}}
    \UnaryInf$\fCenter (!A \land !B) \lif !A$
  \end{prooftree}
  आरंभीच्या क्रमवर्तीची गोडेल संख्या $p_0 = \tuple{0,
    \Gn{!A \Sequent !A}}$ असेल. $\LeftR{\land}$ च्या निष्कर्षावर
    संपणाऱ्या !!{derivation} ची गोडेल संख्या $p_1 =
    \tuple{1, p_0, \Gn{!A \land !B \Sequent !A}, 9}$ असेल. येथे $1$
    हे $\LeftR{\land}$ ला एकच आधारविधान असल्यामुळे आले आहे;
    $!A \land !B \Sequent !A$ या निष्कर्षाची गोडेल संख्या पुढे आहे;
    आणि $9$ हा $\LeftR{\land}$ चा संकेतांक आहे. मग संपूर्ण
    !!{derivation} ची गोडेल संख्या
    $\tuple{1, p_1, \Gn{\Sequent (!A \land !B) \lif !A}, 14}$,
    म्हणजेच
  \[
  \tuple{1, \tuple{1, \tuple{0, \Gn{!A \Sequent !A}}, \Gn{!A \land !B \Sequent !A}, 9},
    \Gn{\Sequent (!A \land !B) \lif !A}, 14}.
  \]
\end{ex}

\begin{explain}
!!{derivation}s चे निरूपण ठरल्यानंतर, अशा !!{derivation}s वर आदिम पुनरावर्ती रीतीने
क्रिया करता येतात आणि त्यांचे आवश्यक गुणधर्म व संबंध तसेच व्यक्त करता
येतात, हेही दाखवावे लागेल. काही क्रिया सोप्या आहेत. उदाहरणार्थ,
!!a{derivation} ची गोडेल संख्या~$p$ दिली असता,
$\fn{EndSequent}(p) = (p)_{(p)_0+1}$ हे तिच्या अंतिम क्रमवर्तीची
गोडेल संख्या देते आणि $\fn{LastRule}(p) = (p)_{(p)_0+2}$ हे तिच्या
अखेरच्या नियमाचा संकेतांक देते. पुढीलप्रमाणे परिभाषित केलेला
क्रमिकेच्या संकेतांकाची वैधता तपासण्यासाठी आदिम पुनरावर्ती अट घ्या:
\[
\fn{Seq}(u) \defiff u=0 \lor u=\prod_{i<\len{u}}p_i^{(u)_i+1}.
\]
ती रिक्त क्रमिकेचे दोन्ही संकेतांक स्वीकारते आणि अतिरिक्त अविभाज्य
घटक असलेले अवैध संकेतांक वगळते.
$\fn{Sequent}(s)$ गुणधर्म
\[
  \begin{gathered}
  s=\tuple{(s)_0,(s)_1} \land \fn{Seq}((s)_0) \land \fn{Seq}((s)_1)\\
  {}\land \bforall{i<\len{(s)_0} + \len{(s)_1}}{\fn{Sent}(((s)_0 \concat (s)_1)_i)}
  \end{gathered}
\]
$s$ साठी खरा असतो, तर आणि तरच $s$ ही !!{sentence}s असलेल्या
क्रमवर्तीची गोडेल संख्या असते. काही इतर क्रिया अधिक कठीण आहेत; त्या कशा कराव्यात याची
किमान रूपरेषा आपण पाहू. $\pi$ ही $\Gamma$ पासून $!A$ ची
!!a{derivation} आहे, हा संबंध $\pi$ आणि $!A$ यांच्या गोडेल संख्यांचा
आदिम पुनरावर्ती संबंध आहे, हे दाखवणे आपले उद्दिष्ट आहे.
\end{explain}

\begin{prop}\ollabel{prop:followsby}
  गोडेल संख्या~$p$ असलेल्या !!{derivation}~$\pi$ मधील अखेरचे निगमन
  योग्य असेल, तर आणि तरच खरा असणारा $\fn{Correct}(p)$ गुणधर्म
  आदिम पुनरावर्ती आहे.
\end{prop}

\begin{proof}
  $\Gamma \Sequent \Delta$ ही आरंभीची क्रमवर्ती दोनपैकी एका
  परिस्थितीत असते: एखादे !!a{sentence}~$!A$ असे असते की
  $\Gamma \Sequent \Delta$ म्हणजे $!A \Sequent !A$ असते;
  किंवा एखादे बंद पद~$t$ असे असते की
  $\Gamma \Sequent \Delta$ म्हणजे $\emptyset \Sequent \eq[t][t]$ असते.
  गोडेल संख्यांच्या भाषेत, मूलभूत आरंभीच्या क्रमवर्तींसाठी
  $\fn{InitialSeq}(s)$ पुढील अट पूर्ण झाली तर आणि तरच खरे असते;
  चिन्हांकित आवृत्त्यांत खाली दिलेल्या अटीही जोडायच्या आहेत:
  \begin{align*}
  \bexists{x < s}{} (\fn{Sent}(x) & \land
  s = \tuple{\tuple{x},\tuple{x}})
  \lor {}\\
  \bexists{t<s}{} (\fn{ClTerm}(t) & \land
  s = \tuple{0, \tuple{\Gn{{\eq}(} \concat t \concat \Gn{,} \concat t \concat \Gn{)}}}).
  \end{align*}
  \iftag{prvTrue}{या चिन्हांकित आवृत्तीत $s=\tuple{0,\tuple{\Gn{\ltrue}}}$
  हीदेखील आरंभीची क्रमवर्ती आहे.}{}
  \iftag{prvFalse}{या चिन्हांकित आवृत्तीत $s=\tuple{\tuple{\Gn{\lfalse}},0}$
  हीदेखील आरंभीची क्रमवर्ती आहे.}{}

  प्रत्येक निगमन नियम~$R$ साठी $\fn{FollowsBy}_R(p)$ हा संबंधही
  आदिम पुनरावर्ती आहे, हे दाखवावे लागेल. $p$ ही !!{derivation}~$\pi$
  ची गोडेल संख्या असेल आणि $\pi$ ची अंतिम क्रमवर्ती तिच्या तात्काळ
  $\pi$ च्या उप-!!{derivation}s च्या अंतिम क्रमवर्तींपासून $R$ च्या योग्य वापराने
  निष्पन्न होत असेल, तर आणि तरच $\fn{FollowsBy}_R(p)$ खरे असते.

  \RightR{\land} नियमाचे उदाहरण सोपे आहे. $\pi$ चा शेवट
  $\RightR{\land}$ च्या योग्य वापराने होत असेल, तर तिचे रूप असे असते:
  \begin{prooftree}
    \AxiomC{}
    \RightLabel{$\pi_1$}
    \Deduce$\Gamma \fCenter \Delta, !A$

    \AxiomC{}
    \RightLabel{$\pi_2$}
    \Deduce$\Gamma \fCenter \Delta, !B$

    \RightLabel{\RightR\land}
    \BinaryInf$\Gamma \fCenter \Delta, !A \land !B$
  \end{prooftree}
  त्यामुळे !!{derivation}~$\pi$ मधील अखेरचे निगमन $\RightR{\land}$ चा योग्य वापर
  असेल, तर आणि तरच $\Gamma$ व $\Delta$ या !!{sentence}s च्या क्रमिका
  आणि $!A$, $!B$ ही दोन !!{sentence}s अशी असतात की $\pi_1$ ची
  अंतिम क्रमवर्ती $\Gamma \Sequent \Delta, !A$, $\pi_2$ ची अंतिम
  क्रमवर्ती $\Gamma \Sequent \Delta, !B$, आणि $\pi$ ची अंतिम
  क्रमवर्ती $\Gamma \Sequent \Delta, !A \land !B$ असते.
  हेच आता गोडेल संख्यांत मांडायचे आहे. $s = \Gn{\Gamma \Sequent \Delta}$
  असेल, तर $(s)_0 = \Gn{\Gamma}$ आणि $(s)_1 = \Gn{\Delta}$.
  म्हणून $\fn{FollowsBy}_{\RightR{\land}}(p)$ पुढील अट पूर्ण झाली
  तर आणि तरच खरे असते:
\begin{align*}
& \bexists{g < p}{\bexists{d < p}{\bexists{a < p}{\bexists{b < p}{\quad}}}} \\
& \qquad \fn{Sequent}(\tuple{g,d}) \land \fn{Sent}(a) \land \fn{Sent}(b) \land {}\\
& \qquad \fn{EndSequent}(p) =
  \tuple{g, d \concat
    \tuple{\Gn{(} \concat a \concat \Gn{\land} \concat b \concat \Gn{)}}} \land {} \\
& \qquad \fn{EndSequent}((p)_1) = \tuple{g, d \concat \tuple{a}} \land {}\\
& \qquad \fn{EndSequent}((p)_2) = \tuple{g, d \concat \tuple{b}} \land {}\\
& \qquad (p)_0 = 2 \land \fn{LastRule}(p) = 10.
\end{align*}
या ओळी अनुक्रमे पुढील गोष्टी व्यक्त करतात: $\Gamma$ ची गोडेल संख्या
$g$, $\Delta$ ची गोडेल संख्या $d$, !!a{formula}~$!A$ ची गोडेल संख्या $a$, आणि
!!a{formula}~$!B$ ची गोडेल संख्या $b$ आहे; $\pi$ ची अंतिम क्रमवर्ती
$\Gamma \Sequent \Delta, !A \land !B$ आहे; $\pi_1$ ची अंतिम क्रमवर्ती
$\Gamma \Sequent \Delta, !A$ आहे; $\pi_2$ ची अंतिम क्रमवर्ती
$\Gamma \Sequent \Delta, !B$ आहे; आणि $\pi$ ला दोन तात्काळ
उपनिष्पत्ती असून तिचा अखेरचा नियम क्रमांक~$10$ असलेला
$\RightR\land$ आहे.

$\pi$ मधील अखेरचे निगमन $\RightR\lexists$ चा योग्य वापर असेल,
तर आणि तरच $\Gamma$ व $\Delta$ या क्रमिका, $!A$ हे !!a{formula},
$x$ हे चल आणि $t$ हे बंद पद अशी असतात की $\pi$ ची अंतिम क्रमवर्ती
$\Gamma \Sequent \Delta, \lexists[x][!A]$ आणि $\pi_1$ ची अंतिम
क्रमवर्ती $\Gamma \Sequent \Delta, \Subst{!A}{t}{x}$ असते.
म्हणून गोडेल संख्यांत $\fn{FollowsBy}_{\RightR\lexists}(p)$ पुढील
अट पूर्ण झाली तर आणि तरच खरे असते:
\begin{align*}
  & \bexists{g<p}{\bexists{d<p}{\bexists{a<p}{\bexists{x<p}{\bexists{t<p}{\quad}}}}}\\
  & \qquad \fn{Sequent}(\tuple{g,d}) \land \fn{Frm}(a) \land \fn{Var}(x) \land {}\\
  & \qquad \fn{ClTerm}(t) \land {}\\
  & \qquad \fn{EndSequent}(p) =
  \tuple{
    g,
    d \concat \tuple{\Gn{\lexists} \concat x \concat a}
  } \land {}\\
  & \qquad \fn{EndSequent}((p)_1) =
  \tuple{
    g,
    d \concat \tuple{\fn{Subst}(a, t, x)}
  } \land {}\\
  & \qquad (p)_0 = 1 \land \fn{LastRule}(p) = 18.
\end{align*}
प्रस्तावित निष्पत्तीगाठीचा संकेतांक नेमक्या पुढीलपैकी एका रूपात असावा:
\begin{align*}
\fn{ProofNode}(p) \defiff {}& p=\tuple{0,\fn{EndSequent}(p)} \lor {}\\
& p=\tuple{1,(p)_1,\fn{EndSequent}(p),\fn{LastRule}(p)} \lor {}\\
& p=\tuple{2,(p)_1,(p)_2,\fn{EndSequent}(p),\fn{LastRule}(p)}.
\end{align*}
ही अट अनावश्यक अतिरिक्त घटक किंवा चुकीची गाठीची स्थानसंख्या वगळते.
मग $\fn{Correct}(p)$ ची व्याख्या अशी करतो:
\begin{multline*}
  \fn{ProofNode}(p) \land \fn{Sequent}(\fn{EndSequent}(p)) \land {}\\
  [(\fn{LastRule}(p) = 1 \land
  \fn{FollowsBy}_{\LeftR\Weakening}(p)) \lor \dots \lor {}\\
  (\fn{LastRule}(p) = 20 \land \fn{FollowsBy}_{\eq}(p)) \lor {}\\
  (p)_0 = 0 \land \fn{InitialSeq}(\fn{EndSequent}(p))]
\end{multline*}
पहिली ओळ गाठीच्या संकेतांकाचे नेमके रूप तपासते आणि $p$ ची अंतिम क्रमवर्ती खरोखर !!{sentence}s चीच
क्रमवर्ती आहे, याची खात्री करते. शेवटची ओळ $p$ मध्ये फक्त आरंभीची
क्रमवर्ती असण्याचा प्रसंग हाताळते.
%READERNOTE{SOL6-A065 — मूळ proof-code चाचण्या decoded fields वरच विसंबतात. येथे sequence codes, नेमका दोन-बाजूंचा sequent code, नेमकी initial/unary/binary proof-node रूपे आणि दाखवलेल्या दोन्ही नियमांचे syntax domains तपासले. लांबी दोन असली तरी extra prime factor असलेल्या 42 चा sequent code म्हणून स्वीकार किंवा initial node ला तिसरा अनावश्यक घटक जोडण्याचा स्वीकार यामुळे थांबतो.}
\end{proof}

\begin{prob}
  \olref[inc][art][plk]{prop:followsby} प्रमाणे पुढील गुणधर्म
  परिभाषित करा:
  \begin{enumerate}
  \item $\fn{FollowsBy}_{\Cut}(p)$,
  \item $\fn{FollowsBy}_{\LeftR\lif}(p)$,
  \item $\fn{FollowsBy}_{\eq}(p)$,
  \item $\fn{FollowsBy}_{\RightR{\lforall}}(p)$.
  \end{enumerate}
  शेवटच्या गुणधर्मासाठी, गोडेल संख्या~$p$ असलेल्या
  !!{derivation} मधील अखेरचे निगमन आयगेन चराची अट पूर्ण करते
  का, हे आदिम पुनरावर्ती रीतीने तपासता येते, हेही दाखवावे लागेल.
  म्हणजेच $\RightR{\lforall}$ चा आयगेन चर~$a$ अंतिम
  क्रमवर्तीत येत नाही.
\end{prob}

\begin{prop}
  \ollabel{prop:deriv}
  $p$ ही योग्य !!{derivation}~$\pi$ ची गोडेल संख्या असेल, तर खरा
  असणारा $\fn{Deriv}(p)$ संबंध आदिम पुनरावर्ती आहे.
\end{prop}

\begin{proof}
  !!^a{derivation}~$\pi$ मधील प्रत्येक निगमन एखाद्या नियमाचा योग्य
  वापर असेल, म्हणजेच तिच्या प्रत्येक उप-!!{derivation}s चा शेवट
  योग्य निगमनाने होत असेल, तर ती योग्य असते. म्हणून
  $\fn{Deriv}(p)$ तर आणि तरच
  \[
  \bforall{i<\len{\fn{SubtreeSeq}(p)}}{\fn{Correct}((\fn{SubtreeSeq}(p))_i)}.
  \]
\end{proof}

\begin{prop}
समजा $\Gamma$ हा !!{sentence}s चा आदिम पुनरावर्ती संच आहे.
मग $\Prf[\Gamma](x, y)$ हा संबंध आदिम पुनरावर्ती आहे. हा संबंध
असे व्यक्त करतो: $x$ हा $\Gamma_0 \Sequent !A$ च्या
!!a{derivation}~$\pi$ चा संकेतांक आहे, येथे काही सांत
$\Gamma_0 \subseteq \Gamma$ आहे, आणि $y$ ही $!A$ ची गोडेल संख्या आहे.
\end{prop}

\begin{proof}
समजा ``$y \in \Gamma$'' हे $R_\Gamma(y)$ या आदिम पुनरावर्ती
विधेयाने दिले आहे. $\Prf[\Gamma](x, y)$ हा संबंध आदिम पुनरावर्ती आहे,
हे दाखवायचे आहे. $y$ ही वाक्य~$!A$ ची गोडेल संख्या असेल
आणि $x$ हा अंतिम क्रमवर्ती $\Gamma_0 \Sequent !A$ असलेल्या
$\Log{LK}$-!!{derivation} चा संकेतांक असेल, तर आणि तरच हा संबंध खरा असतो.

मागील प्रस्तावानुसार, $x$ हा $\Log{LK}$ मधील योग्य
!!{derivation}~$\pi$ चा संकेतांक असेल, तर आणि तरच खरा असणारा
$\fn{Deriv}(x)$ गुणधर्म आदिम पुनरावर्ती आहे. $x$ असा संकेतांक असेल,
तर $\fn{EndSequent}(x)$ हा $\pi$ च्या अंतिम क्रमवर्तीचा संकेतांक आहे.
म्हणून $(\fn{EndSequent}(x))_0$ हा तिच्या डाव्या बाजूचा आणि
$(\fn{EndSequent}(x))_1$ हा उजव्या बाजूचा संकेतांक आहे. त्यामुळे
``$\pi$ च्या अंतिम क्रमवर्तीची उजवी बाजू $!A$ आहे'' हे
$\len{(\fn{EndSequent}(x))_1} = 1 \land ((\fn{EndSequent}(x))_1)_0 =
y$ असे व्यक्त करता येते. $\pi$ च्या अंतिम क्रमवर्तीची डावी बाजू
अर्थातच सांत असते; तिच्यातील प्रत्येक वाक्य $\Gamma$ मध्ये आहे,
हे व्यक्त करणे उरते. म्हणून $\Prf[\Gamma](x, y)$ ची व्याख्या अशी करता येते:
\begin{align*}
\Prf[\Gamma](x, y) \defiff {}&
\fn{Deriv}(x) \land {} \\
& \bforall{i <
  \len{(\fn{EndSequent}(x))_0}}{R_\Gamma(((\fn{EndSequent}(x))_0)_i)} \land {}\\
& \len{(\fn{EndSequent}(x))_1} = 1 \land ((\fn{EndSequent}(x))_1)_0 = y.
\end{align*}
\end{proof}

\end{document}
